\section{Nebula-FHE: DAG/GPU-native scheduling}
\label{sec:nebula-fhe-scheduling}

Nebula-FHE is the DAG/GPU-native scheduler for encrypted computation
in Lux. The F-Chain \texttt{drain\_fhe} service from LP-013 is its
canonical realisation; this section describes the layer Nebula-FHE
exposes to non-F-Chain deployments and to encrypted-coprocessor
workloads outside the wave-tick kernel.

\subsection{Job model}

A Nebula-FHE job is a tuple
\[
  \mathsf{FheJob} = (\mathsf{circuit\_dag\_root},
  \mathsf{ciphertext\_inputs},
  \mathsf{evaluation\_key\_id})
\]
where \texttt{circuit\_dag\_root} is the merkle root of the encrypted
circuit DAG, \texttt{ciphertext\_inputs} is the input vector (LWE
ciphertexts in the canonical wire format), and \texttt{evaluation\_key\_id}
references the bootstrap-key + key-switch-key pair the runtime should
use. The job emits an $\mathsf{FheReceipt} = (\mathsf{output\_root},
\mathsf{kernel\_id}, \mathsf{cost})$ where \texttt{output\_root} is the
merkle root over the produced ciphertexts and \texttt{kernel\_id}
identifies the kernel sequence (CPU vs Lux FHE-CUDA vs Lux FHE-HIP).

\subsection{Scheduling policy}

The DAG planner orders FheJobs by lane affinity:

\begin{itemize}
\item \textbf{Lane-local}: independent ciphertexts (no shared
encrypted state) execute in parallel; this is the throughput sweet spot
of the GPU backends.
\item \textbf{Hot-lane}: shared encrypted state (e.g. an encrypted
ERC-20 \texttt{totalSupply}) routes through a deterministic reducer
under LP-010 \S HotLaneMode. The reducer is itself an FHE circuit; it
runs on the same GPU backend but is single-threaded for the contended
slot.
\end{itemize}

The lane assignment is decided up front by the planner, not by the
GPU at execution time. This keeps GPU dispatch deterministic --- a
property the dual-runtime equivalence theorem
(\S\,\ref{sec:cross-runtime-katology}, Theorem~\ref{thm:fhe-dual-runtime-equivalence})
relies on.

\subsection{F-Chain co-residency}

Inside an F-Chain wave tick, Nebula-FHE jobs run as the
\texttt{drain\_fhe} service alongside \texttt{drain\_exec} (plaintext
EVM), \texttt{drain\_validate} (Block-STM), \texttt{drain\_commit}
(root chain), and \texttt{drain\_cert\_lane} (BLS / Ringtail / ML-DSA
/ FChainTFHE). LP-013 \S Wave-tick co-residency specifies the layout
in detail. The same \texttt{MvccSlot} table backs both encrypted and
plaintext SLOAD; Block-STM sees ciphertext as opaque bytes.

\subsection{Outside F-Chain}

For non-F-Chain deployments (e.g. an external encrypted-coprocessor
service consumed by a confidential-ERC-20 dApp on a non-F-Chain L1),
Nebula-FHE exposes the same job model through
\texttt{github.com/luxfi/fhe-coprocessor}. The coprocessor publishes
$\mathsf{FheReceipt}$ commitments via Warp 2.0 envelopes (LP-021v2,
\texttt{WarpEnvelopeV2}); the destination chain verifies the Pulse
over the receipt root and accepts the encrypted result without
trusting the coprocessor unilaterally. \emph{Status: planned for the
v0.1.0-rc2-coproc milestone.}
